\documentclass[../../master-handout.tex]{subfiles}

\begin{document}
\section{Lecture 3 (Wed, May 20)}
\sidenote{\textbf{Scribe:} Shourrya Guha \\ \textbf{Date:} Wed, May 20, 2026}

\subsection{Serial Manipulators and Configuration Space}
A \textbf{Serial Manipulator} is defined as a series of links interconnected by joints. To understand the state of a manipulator, we look at its configuration:
\begin{itemize}
    \item \textbf{Configuration:} The complete specification of all the points of the manipulator.
    \item \textbf{Configuration Space:} The set of all possible configurations.
    \item \textbf{Degrees of Freedom (DOF):} A manipulator has $n$ DOF if its configuration can be specified by exactly $n$ independent numbers.
\end{itemize}

\textbf{Note:} Knowing the angles of all the joints of a manipulator is equivalent to knowing the full configuration of the manipulator. 

The configuration is typically represented by a \textbf{Vector of Joint Angles} (Configuration Vector):\sidenote{Joint variables are the minimal coordinates used to parameterize the manipulator state.}
Figure~\ref{fig:serial-config} summarizes the configuration vector visually.
\begin{fullwidth}
\[
q = \begin{bmatrix} q_1 \\ q_2 \\ \vdots \\ q_n \end{bmatrix} = \begin{bmatrix} \theta_1 \\ \theta_2 \\ \vdots \\ \theta_n \end{bmatrix}
\]
\end{fullwidth}

\begin{marginfigure}
\centering
\includegraphics[width=\linewidth]{lectures/week-2/lecture-3-images/output.png}
\caption{A serial manipulator configuration example.}
\label{fig:serial-config}
\end{marginfigure}

\subsection{Workspace and Degrees of Freedom}
Objects in 3D space inherently have 6 DOF:
\begin{itemize}
    \item 3 for position (translation)
    \item 3 for orientation (rotation)
\end{itemize}

In contrast, a manipulator can have:
\begin{itemize}
    \item $< 6$ DOF (under-actuated for full spatial control)
    \item $= 6$ DOF (fully actuated)
    \item $> 6$ DOF (redundant manipulator)
\end{itemize}

\paragraph{Workspace Definitions:}
The \textbf{Workspace} is the set of all points reachable by the End Effector (EE).\sidenote{This is a position-space description; orientation requirements may shrink the usable set further.}
Figure~\ref{fig:workspace} shows the reachable region used in the discussion below.

\begin{marginfigure}
\centering
\includegraphics[width=\linewidth]{lectures/week-2/lecture-3-images/workspace.png}
\caption{Workspace illustration.}
\label{fig:workspace}
\end{marginfigure}

\begin{itemize}
    \item \textbf{Dexterous Workspace:} The set of all points reachable with \textit{any} orientation of the end effector.
    \item \textbf{Reachable Workspace:} The set of all reachable points (where the EE might not be able to achieve every orientation). 
\end{itemize}
It follows that the Dexterous Workspace is a subset of the Reachable Workspace.

\subsection{Direct (Forward) Kinematics}
Forward kinematics involves calculating the position and orientation of the end effector relative to the base frame, given the joint parameters. 

\textbf{Convention:} A superscript indicates the \textit{reference frame} (With Respect To), while a subscript indicates the frame being described.

The position vector $p^0$ (point in frame $n$ expressed in frame 0) can be calculated as:\sidenote{The superscript and subscript encode the frame convention, so the expression stays unambiguous.}
\begin{fullwidth}
\[
p^0 = R_n^0 p^n + d_n^0
\]
\end{fullwidth}
Where $d_n^0$ is the displacement vector from frame 0 to frame $n$, and the rotation matrix $R_n^0$ is defined by its column vectors (the principal axes of frame $n$ expressed in frame 0):
\begin{fullwidth}
\[
R_n^0 = \begin{bmatrix} | & | & | \\ x_n^0 & y_n^0 & z_n^0 \\ | & | & | \end{bmatrix} \in \mathbb{R}^{3 \times 3}
\]
\end{fullwidth}

\paragraph{Example 1: Pure Translation and Rotation}
Given the displacement vector $d_n^0$ (origin of frame $n$ with respect to 0) and a point vector $p^n$:
Figure~\ref{fig:ex1} illustrates the frame transform used in this example.

\begin{marginfigure}
\centering
\includegraphics[width=\linewidth]{lectures/week-2/lecture-3-images/ex1.png}
\caption{Pure translation and rotation example.}
\label{fig:ex1}
\end{marginfigure}

\begin{fullwidth}
\[
d_n^0 = O_n^0 = \begin{bmatrix} 1 \\ 1 \\ 0 \end{bmatrix}, \quad p^n = \begin{bmatrix} 1 \\ 1 \\ 0 \end{bmatrix}
\]
\end{fullwidth}
And the rotation matrix $R_n^0$:
\begin{fullwidth}
\[
R_n^0 = \begin{bmatrix} 0 & -1 & 0 \\ 1 & 0 & 0 \\ 0 & 0 & 1 \end{bmatrix}
\]
\end{fullwidth}
We calculate $p^0$ as follows:
\begin{fullwidth}
\[
p^0 = \begin{bmatrix} 0 & -1 & 0 \\ 1 & 0 & 0 \\ 0 & 0 & 1 \end{bmatrix} \begin{bmatrix} 1 \\ 1 \\ 0 \end{bmatrix} + \begin{bmatrix} 1 \\ 1 \\ 0 \end{bmatrix}
= \begin{bmatrix} -1 \\ 1 \\ 0 \end{bmatrix} + \begin{bmatrix} 1 \\ 1 \\ 0 \end{bmatrix} = \begin{bmatrix} 0 \\ 2 \\ 0 \end{bmatrix}
\]
\end{fullwidth}

\paragraph{Example 2: Fundamental Rotations}
Rotation about the Z-axis by angle $\theta$:
Figure~\ref{fig:rotz} shows the reference frame for the $z$-axis rotation.

\begin{marginfigure}
\centering
\includegraphics[width=\linewidth]{lectures/week-2/lecture-3-images/ex2.png}
\caption{Rotation about the $z$-axis.}
\label{fig:rotz}
\end{marginfigure}


\begin{fullwidth}
\[
R_{z,\theta} = \begin{bmatrix} \cos\theta & -\sin\theta & 0 \\ \sin\theta & \cos\theta & 0 \\ 0 & 0 & 1 \end{bmatrix}
\]
\end{fullwidth}
Rotation about the Y-axis by angle $\theta$:
Figure~\ref{fig:roty} shows the analogous $y$-axis rotation setup.

\begin{marginfigure}
\centering
\includegraphics[width=\linewidth]{lectures/week-2/lecture-3-images/ex3.png}
\caption{Rotation about the $y$-axis.}
\label{fig:roty}
\end{marginfigure}

\begin{fullwidth}
\[
R_{y,\theta} = \begin{bmatrix} \cos\theta & 0 & \sin\theta \\ 0 & 1 & 0 \\ -\sin\theta & 0 & \cos\theta \end{bmatrix}
\]
\end{fullwidth}

\subsection{Homogeneous Transformations}
To combine rotation and translation into a single matrix multiplication, we use homogeneous coordinates. The transformation can be written as:
\begin{fullwidth}
\[
P^0 = H_n^0 P^n
\]
\[
\begin{bmatrix} X_0 \\ 1 \end{bmatrix} = \begin{bmatrix} R & d \\ 0 & 1 \end{bmatrix} \begin{bmatrix} X_n \\ 1 \end{bmatrix}
\]
\end{fullwidth}
Where $H \in SE(3)$, which is an element of the Special Euclidean motion group.

\subsection{Properties of Rotation Matrices}
A valid rotation matrix belongs to the Special Orthogonal Group, denoted as $R \in SO(3)$:
\begin{fullwidth}
\[
R = \begin{bmatrix} | & | & | \\ x & y & z \\ | & | & | \end{bmatrix}
\]
\end{fullwidth}
Key properties include:
\begin{enumerate}
    \item \textbf{Unit Norm:} The column vectors are unit vectors, so $||x|| = ||y|| = ||z|| = 1$.
    \item \textbf{Orthogonality:} The column vectors are mutually orthogonal, meaning their dot products are zero: $x^T y = 0, \quad y^T z = 0, \quad z^T x = 0$.
    \item \textbf{Determinant:} The determinant of a rotation matrix is strictly positive one: $\det(R) = +1$.
\end{enumerate}

\end{document}